% VI42-DETAILED-EXERCISE-SOLUTIONS
% VI42-DETAILED-EXERCISE-01
\subsection*{Exercise VI.42.1}
\begin{solution}
A trivialization \(\mathcal L|_U\simeq\mathcal O_U\) sends \(1\) to a section \(e\) that generates every stalk on \(U\).  Thus \(e\) is a local basis of one element, so \(\mathcal L\) is locally free of rank one.  Conversely, a rank-one local basis \(e\) defines the isomorphism \(a\mapsto ae\) from \(\mathcal O_U\) to \(\mathcal L|_U\).  Hence the two definitions are literally equivalent, not merely analogous.
\end{solution}

% VI42-DETAILED-EXERCISE-02
\subsection*{Exercise VI.42.2}
\begin{solution}
On \(U_i\cap U_j\cap U_k\),
\[
e_i=g_{ij}e_j=g_{ij}g_{jk}e_k.
\]
But by definition \(e_i=g_{ik}e_k\).  Since \(e_k\) is a basis, coefficients are unique, so \(g_{ij}g_{jk}=g_{ik}\).  Taking \(k=i\) also recovers \(g_{ij}g_{ji}=1\).
\end{solution}

% VI42-DETAILED-EXERCISE-03
\subsection*{Exercise VI.42.3}
\begin{solution}
If \(e_i'=u_ie_i\) and \(e_j'=u_je_j\), then
\[
e_i'=u_i g_{ij}e_j=u_ig_{ij}u_j^{-1}e_j'.
\]
Therefore \(g'_{ij}=u_i g_{ij}u_j^{-1}\).  This is the multiplicative coboundary action; it explains why changing local frames alters a cocycle but not the isomorphism class of the line bundle.
\end{solution}

% VI42-DETAILED-EXERCISE-04
\subsection*{Exercise VI.42.4}
\begin{solution}
Choose an open cover on which both \(\mathcal L\) and \(\mathcal M\) are trivial.  On such a \(U\),
\[
(\mathcal L\otimes\mathcal M)|_U\simeq\mathcal O_U\otimes_{\mathcal O_U}\mathcal O_U\simeq\mathcal O_U.
\]
Thus the tensor product is locally free of rank one.  If the transition functions are \(g_{ij}\) and \(h_{ij}\), the tensor product has transitions \(g_{ij}h_{ij}\).
\end{solution}

% VI42-DETAILED-EXERCISE-05
\subsection*{Exercise VI.42.5}
\begin{solution}
On a trivializing open choose a frame \(e\) and dual frame \(e^\vee\) with \(e^\vee(e)=1\).  Evaluation gives
\[
e\otimes e^\vee\longmapsto1,
\]
hence an isomorphism \(\mathcal L\otimes\mathcal L^\vee\simeq\mathcal O_X\) locally.  These maps are canonical and therefore agree on overlaps, so they glue globally.
\end{solution}

% VI42-DETAILED-EXERCISE-06
\subsection*{Exercise VI.42.6}
\begin{solution}
If \(\mathcal L\) has transitions \(g_{ij}\), then \(\mathcal L^\vee\) has transitions \(g_{ij}^{-1}\).  Hence
\[
\mathcal L\otimes\mathcal L^\vee
\]
has transition functions \(g_{ij}g_{ij}^{-1}=1\), so it is trivial.  Consequently \([\mathcal L^\vee]=-[\mathcal L]\) in \(\mathrm{Pic}(X)\).
\end{solution}

% VI42-DETAILED-EXERCISE-07
\subsection*{Exercise VI.42.7}
\begin{solution}
Pullback preserves local freeness and satisfies the canonical tensor identity
\[
f^*(\mathcal L\otimes\mathcal M)\simeq f^*\mathcal L\otimes f^*\mathcal M.
\]
Therefore \([\,\mathcal L\,]\mapsto[\,f^*\mathcal L\,]\) preserves addition in the Picard group.  Also \(f^*\mathcal O_X\simeq\mathcal O_Y\), so the identity class is preserved.
\end{solution}

% VI42-DETAILED-EXERCISE-08
\subsection*{Exercise VI.42.8}
\begin{solution}
For \(D=(U_i,f_i)\), the standard frame of \(\mathcal O_X(D)|_{U_i}=f_i^{-1}\mathcal O_{U_i}\) is \(e_i=f_i^{-1}\).  On \(U_i\cap U_j\),
\[
e_i=f_i^{-1}=\frac{f_j}{f_i}f_j^{-1}=\frac{f_j}{f_i}e_j.
\]
Thus the transition function is \(f_j/f_i\), which is a regular unit precisely because the \(f_i\) define a Cartier divisor.
\end{solution}

% VI42-DETAILED-EXERCISE-09
\subsection*{Exercise VI.42.9}
\begin{solution}
If \(D=\operatorname{div}(f)\), then \(\mathcal O_X(D)=f^{-1}\mathcal O_X\subset K(X)\).  Multiplication by \(f\) is an \(\mathcal O_X\)-linear isomorphism
\[
f^{-1}\mathcal O_X\xrightarrow{\sim}\mathcal O_X.
\]
This proves that principal Cartier divisors map to the identity in the Picard group.
\end{solution}

% VI42-DETAILED-EXERCISE-10
\subsection*{Exercise VI.42.10}
\begin{solution}
If \(D\) and \(E\) have local equations \(f_i\) and \(g_i\), then \(D+E\) has equation \(f_ig_i\).  Its local frame is \((f_ig_i)^{-1}\), while the tensor product frame is \(f_i^{-1}\otimes g_i^{-1}\).  Multiplication in the rational function field identifies these frames and gives
\[
\mathcal O_X(D+E)\simeq\mathcal O_X(D)\otimes\mathcal O_X(E).
\]
\end{solution}

% VI42-DETAILED-EXERCISE-11
\subsection*{Exercise VI.42.11}
\begin{solution}
At the generic point \(\eta\), the stalk \(\mathcal L_\eta\) is a one-dimensional vector space over \(K(X)=\mathcal O_{X,\eta}\).  Choose any nonzero \(s_\eta\in\mathcal L_\eta\).  After shrinking around \(\eta\), it is represented by a section, and its compatible germs define a nonzero rational section.  This rational trivialization is the input used to recover a Cartier divisor.
\end{solution}

% VI42-DETAILED-EXERCISE-12
\subsection*{Exercise VI.42.12}
\begin{solution}
Write \(s=f_i e_i\) in local frames.  Then \(fs=(ff_i)e_i\), so the local equations of the new divisor are \(ff_i\).  Therefore
\[
\operatorname{div}(fs)=\operatorname{div}(s)+\operatorname{div}(f).
\]
Thus a different rational trivialization changes the divisor only by a principal divisor, exactly why the line bundle determines a Cartier class rather than a unique divisor.
\end{solution}

% VI42-DETAILED-EXERCISE-13
\subsection*{Exercise VI.42.13}
\begin{solution}
For \(X=\operatorname{Spec}A\) with \(A\) local, an invertible sheaf corresponds to a finitely generated projective \(A\)-module of rank one.  Every finitely generated projective module over a local ring is free, so it is isomorphic to \(A\).  Hence every line bundle is trivial and \(\mathrm{Pic}(X)=0\).
\end{solution}

% VI42-DETAILED-EXERCISE-14
\subsection*{Exercise VI.42.14}
\begin{solution}
Since \(k[t]\) is a PID, every invertible fractional ideal is principal and every rank-one projective module is free.  Equivalently, \(k[t]\) is a UFD and its affine divisor class group vanishes.  Therefore every line bundle on \(\mathbb A^1_k\) is trivial:
\[
\mathrm{Pic}(\mathbb A^1_k)=0.
\]
\end{solution}

% VI42-DETAILED-EXERCISE-15
\subsection*{Exercise VI.42.15}
\begin{solution}
Write \(\mathcal O_{\mathbb P^1}(1)\simeq\mathcal O_{\mathbb P^1}([\infty])\).  On \(U=\mathbb P^1\setminus\{\infty\}\), the Cartier divisor \([\infty]\) restricts to the zero divisor, so
\[
\mathcal O_{\mathbb P^1}(1)|_U\simeq\mathcal O_U.
\]
The same follows from the nowhere-vanishing frame supplied by the homogeneous coordinate nonzero on this chart.
\end{solution}

% VI42-DETAILED-EXERCISE-16
\subsection*{Exercise VI.42.16}
\begin{solution}
Let \(H\) be a hyperplane.  Since \(\mathcal O(d)\simeq\mathcal O(dH)\), the divisor tensor law gives
\[
\mathcal O(a)\otimes\mathcal O(b)
\simeq\mathcal O(aH+bH)
=\mathcal O((a+b)H)
\simeq\mathcal O(a+b).
\]
This works for negative integers as well by taking duals.
\end{solution}

% VI42-DETAILED-EXERCISE-17
\subsection*{Exercise VI.42.17}
\begin{solution}
Projective space is regular, hence locally factorial, so the natural map \(\mathrm{Pic}(\mathbb P^1)\to\mathrm{Cl}(\mathbb P^1)\) is an isomorphism.  VI/41 gives \(\mathrm{Cl}(\mathbb P^1)\simeq\mathbb Z\), generated by a point/hyperplane.  The corresponding line bundle is \(\mathcal O(1)\), so \(\mathrm{Pic}(\mathbb P^1)\simeq\mathbb Z\).
\end{solution}

% VI42-DETAILED-EXERCISE-18
\subsection*{Exercise VI.42.18}
\begin{solution}
Factor \(F\) into irreducible homogeneous factors with multiplicities.  Its zero divisor has total class \(dH\), because the sum of factor degrees is \(d\).  Thus
\[
[V_+(F)]=d[H]
\quad\Longrightarrow\quad
\mathcal O_{\mathbb P^n}(V_+(F))\simeq\mathcal O_{\mathbb P^n}(d).
\]
The statement includes nonreduced multiplicities automatically.
\end{solution}

% VI42-DETAILED-EXERCISE-19
\subsection*{Exercise VI.42.19}
\begin{solution}
Let \(L=V(x,z)\) be the ruling whose class generates \(\mathrm{Cl}(X)\simeq\mathbb Z/2\).  If its class came from a line bundle, it would have a Cartier representative \(D\).  Then \(L-D\) would be principal, hence Cartier, and therefore \(L\) itself would be Cartier.  VI/40 proves that \(L\) is not Cartier at the vertex.  So the nonzero class is outside the image of \(\mathrm{Pic}(X)\).
\end{solution}

% VI42-DETAILED-EXERCISE-20
\subsection*{Exercise VI.42.20}
\begin{solution}
The punctured cone \(U=X\setminus\{o\}\) is regular, hence locally factorial.  Therefore
\[
\mathrm{Pic}(U)\simeq\mathrm{Cl}(U).
\]
Removing a codimension-two point does not change the Weil class group, so VI/41 gives \(\mathrm{Cl}(U)\simeq\mathbb Z/2\).  Consequently
\[
\mathrm{Pic}(U)\simeq\mathbb Z/2.
\]
\end{solution}

% VI42-DETAILED-EXERCISE-21
\subsection*{Exercise VI.42.21}
\begin{solution}
Write \(s=f_ie_i\) in local frames.  On overlaps \(e_i=g_{ij}e_j\), hence \(f_j=f_i g_{ij}\), so the \(f_i\) differ by units and define one Cartier divisor.  Because \(X\) is integral and the section is nonzero generically, each nonzero \(f_i\) is a non-zero-divisor.  Thus the zero scheme is an effective Cartier divisor and \(\mathcal L\simeq\mathcal O_X((s)_0)\).
\end{solution}

% VI42-DETAILED-EXERCISE-22
\subsection*{Exercise VI.42.22}
\begin{solution}
If \(e'=ue\) is another frame, then \(s_i=f_ie=(f_i u^{-1})e'\).  Thus all homogeneous coordinates are multiplied by the same invertible function \(u^{-1}\):
\[
[f_0:\cdots:f_n]=[f_0u^{-1}:\cdots:f_nu^{-1}].
\]
Hence the projective point is independent of the chosen frame, allowing the local maps to glue.
\end{solution}

% VI42-DETAILED-EXERCISE-23
\subsection*{Exercise VI.42.23}
\begin{solution}
For \(f\in K(C)^\times\), the rational function defines a finite morphism \(C\to\mathbb P^1\) when nonconstant.  The fibers over \(0\) and \(\infty\) have the same weighted degree, so the total order of zeros equals the total order of poles.  Therefore
\[
\deg\operatorname{div}(f)=0.
\]
Hence degree descends from divisors to the Picard group.
\end{solution}

% VI42-DETAILED-EXERCISE-24
\subsection*{Exercise VI.42.24}
\begin{solution}
Every line bundle on \(\mathbb P^1\) is \(\mathcal O(d)\) for a unique \(d\), and its degree is \(d\).  Thus
\[
\deg:\mathrm{Pic}(\mathbb P^1)\xrightarrow{\sim}\mathbb Z.
\]
Its kernel is therefore zero, so \(\mathrm{Pic}^0(\mathbb P^1)=0\).
\end{solution}

